%! Author = lili
%! Date = 8/20/26

\newglossaryentry{provirus}{
    type=defi,
    name={Provirus},
    description={In das Wirtsgenom integrierte doppelsträngige DNA-Kopie eines
    Retrovirus, flankiert von zwei LTRs. Sie wird wie ein zelluläres Gen von
    der Wirts-RNA-Polymerase transkribiert}
}


\newglossaryentry{polypyrimidin}{
    first={Polypyrimidin-Folge (polypyrimidine tract, PPT)},
    name={Polypyrimidin-Folge},
    description={Pyrimidin-reiche Sequenz zwischen Verzweigungsstelle und 3'-Spleißstelle. Bindestelle für U2AF2 und
    wichtig für die Erkennung der 3'-Spleißstelle.},
    type=defi
}

\newglossaryentry{proto_onkogen}{
    name={Proto-Onkogen},
    description={todo},
    type=defi
}

\newglossaryentry{wachstumsfaktor}{
    name={Wachstumsfaktor},
    description={todo},
    type=defi
}


\newglossaryentry{verzweigungsstelle}{
    name={Verzweigungsstelle (Branchpoint)},
    description={Konserviertes Adenosin innerhalb des Introns (Konsensus u.\,a. UACUAAC), dessen 2'-OH-Gruppe im ersten Schritt des Spleißens das 5'-Spleißende angreift und so die Lariat-Struktur ausbildet.},
    type=defi
}
\newglossaryentry{comparative-genomics}{
    name={Vergleich von verwandten Genomen},
    description={Wissenschaftsgebiet, das Ähnlichkeiten und Unterschiede zwischen
    DNA-Sequenzen, Genen, der Genreihenfolge, regulatorischen Sequenzen
    und anderen genomischen Merkmalen untersucht, um festzustellen, wie Organismen
    miteinander verwandt sind.},
    type=defi
}


\newglossaryentry{transesterifizierung}{
    name={Transesterifizierung},
    description={Umesterungsreaktion, die dem nukleären Spleißen zugrunde liegt. Zwei aufeinanderfolgende Transesterifizierungen (Angriff des Verzweigungs-2'-OH, dann des Exon-1-3'-OH) trennen das Intron heraus und verknüpfen die Exons.},
    type=defi
}

\newglossaryentry{primer}{
    type=defi,
    name={tRNA-Primer},
    description={Zelluläre tRNA, die sich an die Primerbindestelle der
    retroviralen RNA anlagert und der Reversen Transkriptase das freie
    3$'$-OH-Ende für den Start der Minusstrang-Synthese liefert. Sie wird vor
    der Plusstrang-Synthese wieder entfernt}
}

\newglossaryentry{translokation}{
    name={Translokation},
    description={todo},
    type=defi
}

\newglossaryentry{transspleissen}{
    name={Trans-Spleißen},
    description={Verknüpfung von Exons aus zwei getrennten RNA-Molekülen (im Gegensatz zum normalen \emph{cis}-Spleißen). Prominent bei Trypanosomen (Spliced-Leader-/Mini-Exon-Anhängung).},
    type=defi
}


\newglossaryentry{tdna}{
    name={T-DNA},
    description={todo},
    type=defi
}

\newglossaryentry{ti_plasmid}{
    name={Ti-Plasmid},
    first={Ti-Plasmid (\textbf{T}umor \textbf{i}nducing \textbf{plasmid})},
    description={todo},
    type=defi
}


\newglossaryentry{gametophyt}{
    type=defi,
    name={Gametophyt},
    description={Haploide (1n) Generation im Entwicklungszyklus der Pflanzen. Bei Angiospermen stark reduziert: der
    männliche Gametophyt ist das Pollenkorn (generative Zelle + Pollenschlauchzelle), der weibliche Gametophyt der Embryosack (Eizelle, 2 Synergiden, 2 Polkerne, 3 Antipodenzellen)}
}

\newglossaryentry{telophase}{
    name={Telophase},
    description={die Chromatiden erreichen die entgegen-
    gesetzten Pole, Rückbildung der Kernhülle,
    Dekondensierung der Chromosomen,
    Nucleoli wieder sichtbar},
    type=defi
}
\newglossaryentry{spacer}{
    name={Spacer},
    plural={Spacer},
    description={Nicht-kodierende Sequenzabschnitte in Transkriptionseinheiten. Man unterscheidet transkribierte Spacer (werden mit-transkribiert und später entfernt) und nicht-transkribierte Spacer (enthalten z.\,B. Spacer-Promotoren und Enhancer), besonders relevant bei der rDNA/rRNA.},
    type=defi
}

\newglossaryentry{sporophyt}{type=defi,
    name={Sporophyt},
    description={Diploide (2n) vielzellige Generation; entsteht aus der Zygote und bildet durch Meiose Sporen. Bei
    Blptenpflanzen die dominante Generation (``die Pflanze'')}
}
\newglossaryentry{somatsegregation}{
    name={Somatische Segregation},
    description={Zufaellige Aufteilung unterschiedlicher Organellengenome auf Tochterzellen bei der Zellteilung (relevant bei Heteroplasmie, v.a. in Pflanzen), sodass verschiedene Gewebe unterschiedliche Organellengenotypen tragen koennen.},
    type=defi
}

\newglossaryentry{selbstspleissen}{
    name={Selbst-Spleißen (self-splicing)},
    description={Spleißen ohne Spleißosom, bei dem das Intron als Ribozym selbst die Reaktion katalysiert. Vorkommen u.\,a.\ in Tetrahymena, Mitochondrien, Chloroplasten, Schleimpilzen, Bacteriophage T4.},
    type=defi
}


\newglossaryentry{replikative_transposition}{
    name={replikative Transposition},
    description={Transposition, bei der der Donor unverändert bleibt und der Rezipient eine Kopie des Transposons erhält},
    type=defi
}

\newglossaryentry{reversetranskriptase}{
    type=defi,
    name={Reverse Transkriptase},
    description={RNA-abhängige DNA-Polymerase der Retroviren und
    Retrotransposons. Synthetisiert ausgehend von einem tRNA-Primer den
    Minusstrang der DNA an der retroviralen RNA und besitzt zusätzlich eine
    RNase\,H-Aktivität, die den RNA-Anteil des Hybrids abbaut}
}


\newglossaryentry{rnaseh}{
    type=defi,
    name={RNase H},
    description={Enzymaktivität der Reversen Transkriptase, die den RNA-Strang
    eines RNA-DNA-Hybrids abbaut. Sie legt das 3$'$-Ende der neu synthetisierten
    DNA frei und erzeugt die RNA-Fragmente, die später als Primer der
    Plusstrang-Synthese dienen}
}

\newglossaryentry{rnaediting}{
    name={RNA-Editierung (RNA editing)},
    description={Post-transkriptionale Änderung der RNA-Sequenz (Basenumwandlung, U-Insertion/-Deletion), sodass die
    reife mRNA von der genomischen Sequenz abweicht. Häufig in Mitochondrien, Plastiden und Trypanosomen, aber auch in
    ``normalen'' Säugetier-Transkripten.},
    type=defi
}

\newglossaryentry{rflp}{
    first={Restriktionsfragmentlängenpolymorphismus (RFLP)},
    name={RFLP},
    description={Die Tatsache, dass genomische DNA verschiedener Individuen unterschiedliche Schnittstellen für Restriktionsenzyme aufweist.},
    type=defi
}
% \cite{docflexRFLP}

\newglossaryentry{ras}{
    name={RAS},
    description={todo},
    type=defi
}

\newglossaryentry{r_bande}{
    name={R-Bande},
    description={offene, früh replizierende und genreiche Bereiche},
    type=defi
}

\newglossaryentry{q_g_bande}{
    name={Q-/G-Bande},
    description={kondensierte, spät replizierende und genarme Bereiche},
    type=defi
}

\newglossaryentry{prophase}{
    name={Prophase},
    description={Kondensierung des Chromatins,
    Chromosomen werden sichtbar,
    Spindelapparat bildet sich},
    type=defi
}

\newglossaryentry{prometaphase}{
    name={Prometaphase},
    description={Kernhülle zerfällt, Kinetochore,
    Spindelapparat vollständig},
    type=defi
}

\newglossaryentry{onkogen}{
    name={Onkogen},
    description={Modifizierte Varianten von im normalen Genom vorhandenen Genen},
    type=defi
}

\newglossaryentry{macrophage}{
    name={Macrophage},
    description={todo},
    type=defi
}

\newglossaryentry{metaphase}{
    name={Metaphase},
    description={Bildung der Metaphasenplatte:
    Chromosomen ordnen sich in einer Ebene
    zwischen den Spindelpolen},
    type=defi
}

\newglossaryentry{reassoziation}{
    name={Reassoziationskinetik},
    description={Geschwindigkeiten, mit denen sich komplementäre DNA-Stränge zusammenlagern},
    type=defi
}

\newglossaryentry{tcr}{
    name={TCR},
    description={todo},
    type=defi
}

\newglossaryentry{shm}{
    name={somatische Hypermutation},
    description={todo},
    type=defi
}


\newglossaryentry{schwere_kette}{
    name={schwere kette},
    description={todo},
    type=defi
}

\newglossaryentry{cregion}{
    name={cregion},
    description={todo},
    type=defi
}

\newglossaryentry{vregion}{
    name={vregion},
    description={todo},
    type=defi
}


\newglossaryentry{hapten}{
    name={hapten},
    description={todo},
    type=defi
}


\newglossaryentry{antigene_determinante}{
    name={antigene determinante},
    description={todo},
    type=defi
}

\newglossaryentry{gewebeverträglichkeitskomplex}{
    name={Gewebeverträglichkeitskomplex},
    description={todo},
    type=defi
}

\newglossaryentry{tlr}{
    name={TLR},
    description={todo},
    type=defi
}

\newglossaryentry{mamp}{
    name={MAMP},
    description={todo},
    type=defi
}

\newglossaryentry{prr}{
    name={PRR},
    description={todo},
    type=defi
}

\newglossaryentry{mhc}{
    first={MHC-Locus / major histo-compatibility complex},
    name={major histo-compatibility complex},
    description={todo},
    type=defi
}

\newglossaryentry{lariat}{
    name={Lariat-Struktur (Lasso)},
    description={Lasso-förmige Struktur des herausgespleißten Introns, entsteht durch eine 5'{}--2'-Phosphodiesterbindung zwischen dem 5'-G des Introns und dem 2'-OH des Verzweigungs-Adenosins.},
    type=defi
}

\newglossaryentry{lymphocyte_dna}{
    name={lymphocyte dna},
    description={todo},
    type=defi
}


\newglossaryentry{leichte_kette}{
    name={leichte kette},
    description={todo},
    type=defi
}

\newglossaryentry{ltr}{
    type=defi,
    name={LTR (Long Terminal Repeat)},
    description={Identische, aus U3--R--U5 aufgebaute Sequenzwiederholung an
    beiden Enden der integrierten proviralen DNA. Die LTRs sind im RNA-Genom
    nicht vorhanden, sondern entstehen erst durch den Mechanismus der reversen
    Transkription mit ihren zwei Sprüngen; sie enthalten Promotor- und
    Enhancer-Elemente}
}

\newglossaryentry{ltrretrotransposon}{
    type=defi,
    name={LTR-Retrotransposon},
    description={Retroposon mit langen terminalen Repeats, Reverser Transkriptase
    und/oder Integrase, das Introns enthalten kann (entfernt in subgenomischer
    mRNA) und 4--6\,bp Zielsequenzverdopplung erzeugt. Beispiele: Ty in
    \emph{S.\,cerevisiae}, copia in \emph{D.\,melanogaster}. Unterscheidet sich
    vom Retrovirus im Wesentlichen durch das fehlende \emph{env}-Gen}
}

\newglossaryentry{transposon_klasse_i}{
    name={Klasse I Transposon},
    description={Mobile RNA Zwischenstufe, deren Vermehrung abhängig von der Herstellung von DNA-Kopien der RNA-
    Transkripte, DNA integriert an neuen Stellen im Genom ist; auch: Retrotransposon},
    type=defi
}

\newglossaryentry{transposon_klasse_ii}{
    name={Klasse II Transposon},
    description={Intermediäres DNA Element mit direkter Transposition und Vermehrung im Genom; auch: Retroposon},
    type=defi
}

\newglossaryentry{konservative_transposition}{
    name={konservative Transposition},
    description={Transposition, bei der der Donor eine Bruchstelle an der (ehemaligen) Stelle des Transposons davon
    trägt und der Rezipient das Transposon erhält},
    type=defi
}

\newglossaryentry{kopplungskarte}{
    name={Kopplungskarte (genetische Karte)},
    description={Zeigt Abstaende zwischen Loci in Einheiten, die auf Rekombinationshaeufigkeiten beruhen. Begrenzt durch die Notwendigkeit, Rekombinationen zwischen variablen Markern zu beobachten. Die Rekombinationshaeufigkeit ist proportional zum physikalischen Abstand der Loci.},
    type=defi
}
\newglossaryentry{innate_immunity}{
    name={innate immunity},
    description={todo},
    type=defi
}

\newglossaryentry{ideogramm}{
    name={Idiogramm},
    description={Schematische Darstellung von Chromosomen, ausgerichtet am Centromer},
    type=defi
}


\newglossaryentry{interphase}{
    name={Interphase},
    description={todo},
    type=defi
}


\newglossaryentry{inversion}{
    name={Inversion},
    description={todo},
    type=defi
}

\newglossaryentry{integrase}{
    type=defi,
    name={Integrase},
    description={Vom \emph{pol}-Bereich codiertes Enzym, das die provirale DNA in
    das Wirtsgenom einbaut. Es erzeugt an den LTRs um zwei Basen
    zurückgesetzte 3$'$-Enden, schneidet die Ziel-DNA versetzt und verknüpft
    die LTR-3$'$-Enden mit den 5$'$-Enden des Targets}
}

\newglossaryentry{introndefinition}{
    name={Intron-Ekennungsmodus},
    description={Erkennungsmodus, bei dem die zusammengehörigen Spleißstellen über das (kurze) Intron hinweg definiert
    werden (``interactions across the intron''). Typisch für kurze Introns.},
    type=defi
}

\newglossaryentry{hklasse}{
    name={H-Klasse},
    description={todo},
    type=defi
}

\newglossaryentry{hla}{
    name={HLA-Komplex},
    description={todo},
    type=defi
}

\newglossaryentry{hämophilie_a}{
    name={Hämophilie A},
    description={Krankheit aufgrund von Mangel an Gerinnungsfaktor VIII. Folge: Auftreten spontaner, unkontrollierbarer
    innerer Blutungen. Hat einen geschlechtsgebundenen Erbgang},
    type=defi
}

\newglossaryentry{hämophilie_b}{
    name={Hämophilie B},
    description={Krankheit aufgrund von Mangel an Blutgerinnungsfaktor IX},
    type=defi
}

\newglossaryentry{hemi_methylierung}{
    name={Hemi-Methylierung},
    description={todo},
    type=defi
}

\newglossaryentry{gtag}{
    name={GT-AG-Regel (Jambon-Regel)},
    description={Regel, nach der Introns fast immer mit GT (im Transkript GU) beginnen und mit AG enden. Auf RNA-Ebene
    korrekter ``GU-AG-Regel''. Gilt für die ``major''/U2-Typ-Introns (kanonische Spleißstellen, ca.\ 98\,\%).},
    type=defi
}

\newglossaryentry{extranukleaer}{
    name={Extranukläere Gene},
    description={Gene ausserhalb des Zellkerns (in Mitochondrium oder Chloroplast). Sie werden in demselben Organellenkompartiment transkribiert und translatiert, in dem sie liegen.},
    type=defi
}

\newglossaryentry{dpe}{
    name={Downstream Promoter Element (DPE)},
    plural={Downstream Promoter Elements (DPE)},
    description={Core-Promotorelement der RNA-Polymerase II, das stromabwärts des Transkriptionsstarts liegt und insbesondere bei TATA-losen Promotoren zur Initiation der Transkription beiträgt.},
    type=defi
}


\newglossaryentry{exondefinition}{
    name={Exon-Ekennungsmodus},
    description={Erkennungsmodus, bei dem die Spleißfaktoren zunächst über das Exon hinweg interagieren (``interactions
    across the exon'') und erst danach zum Intron-überspannenden Komplex umschalten. Wichtig bei langen Introns.},
    type=defi
}


\newglossaryentry{exchange_factor}{
    name={exchange factor},
    description={todo},
    type=defi
}


\newglossaryentry{epitope}{
    name={epitope},
    description={todo},
    type=defi
}



\newglossaryentry{deaminase}{
    name={Deaminase},
    description={Enzym, das eine Aminogruppe abspaltet. Bei der RNA-Editierung: Cytidin-Deaminase (C\,$\to$\,U) und Adenosin-Deaminase (A\,$\to$\,I, Inosin).},
    type=defi
}


\newglossaryentry{thalassaemie}{
    name={$\beta$-Thalassämie},
description={Humane Erbkrankheit mit Hämoglobin-Mangel. Kann durch Spleiß-Mutationen entstehen, die einen ``falschen
Spleißort'' aktivieren und so das $\beta$-Globin-Transkript fehlprozessieren},
type=defi
}


\newglossaryentry{contig}{
name={Contig},
description={Aus ueberlappenden Klonen bzw. Sequenz-Reads zusammengesetzter, zusammenhaengender ("contiguous") DNA-Bereich.},
type=defi
}


\newglossaryentry{cwertparadox}{
name={C-Wert-Paradoxon},
description={Der C-Wert ist die haploide Genomgroesse (DNA-Menge) einer Art. Paradox: die Genomgroesse korreliert bei Eukaryoten nicht mit der Komplexitaet oder Genzahl des Organismus (grosse Genome enthalten viel nicht-kodierende/repetitive DNA).},
type=defi
}


\newglossaryentry{altspleissen}{
name={Alternatives Spleissen},
description={Exons können unterschiedlich kombiniert werden, sodass aus einem Gen mehrere verschiedene
mRNAs/Proteine entstehen. Ein Grund, warum "ein Gen = ein Protein" für Eukaryoten nicht zutrifft.},
type=defi
}
\newglossaryentry{annotation}{
name={Annotation},
description={Das Zuweisen von Struktur und Funktion zu einer DNA-Sequenz: Auffinden von Genen, Exon-Intron-Struktur, kodierenden Bereichen und funktionelle Charakterisierung.},
type=defi
}

\newglossaryentry{anaphase}{
name={Anaphase},
description={gepaarte Kinetochore trennen sich,
Chromatiden wandern zu den entgegengesetzten Polen},
type=defi
}

